When does a technology transition from luxury to necessity? For broadband internet, this question has profound implications: if demand becomes price-inelastic, subsidies and price regulations---the dominant policy tools of the past two decades---lose their effectiveness. The transformation of broadband from discretionary communication service to essential infrastructure for economic participation, education, healthcare, and civic engagement \citep{katz2010impact, bertschek2016drivers} suggests this transition may already be underway. Yet empirical evidence on whether and when broadband demand became price-insensitive remains surprisingly scarce, despite substantial public investments in broadband infrastructure worldwide \citep{oecd2020broadband}.

This paper provides such evidence by examining how broadband price elasticity evolved over 2010--2024 in 33 countries: 27 European Union (EU) member states and 6 Eastern Partnership (EaP) countries. We leverage substantial cross-country and temporal variation in broadband prices and adoption rates to estimate time-varying elasticities using two-way fixed effects models. Our analysis reveals a striking structural transformation: Eastern Partnership countries exhibited highly price-elastic demand ($\varepsilon = \EaPBaseBr$) during 2010--2019, but transitioned to price-inelastic demand ($\varepsilon \approx 0$) by 2020--2024, indicating broadband's evolution from discretionary service to essential necessity.

\subsection{Research Questions and Motivation}

Three core questions motivate this research. First, \textit{how has broadband demand elasticity evolved over the past decade and a half?} Existing literature predominantly estimates elasticities at single points in time \citep{dauvin2014estimating, madden2015demand}, yet technological change and digital economy expansion suggest elasticity may be time-varying as services transition from luxury to necessity \citep{hausman2001private}. 

Second, \textit{what explains regional heterogeneity in price sensitivity?} Eastern Partnership countries---comprising Armenia, Azerbaijan, Belarus, Georgia, Moldova, and Ukraine---represent lower-income markets where affordability concerns might generate higher price elasticity compared to wealthier EU markets. Understanding this heterogeneity is crucial for tailoring policy interventions to different development contexts \citep{waverman2001telecommunications}.

Third, \textit{did COVID-19 fundamentally alter broadband demand?} The pandemic forced abrupt transitions to remote work, online education, and digital service delivery \citep{oecd2021covid}, potentially transforming broadband from convenience to necessity overnight. However, distinguishing COVID's causal impact from pre-existing trends requires careful econometric identification, which we address through placebo tests.

\subsection{Contributions}

This paper makes four principal contributions to telecommunications economics and policy.

\textbf{First}, we document time-varying demand elasticity across a 15-year panel, revealing broadband's structural transformation. Unlike prior single-period estimates, our year-by-year analysis shows elasticity declined gradually from 2015 onward, not suddenly in 2020. This challenges the narrative of COVID-19 as an exogenous shock and instead suggests a decade-long evolution driven by digital economy expansion.

\textbf{Second}, we demonstrate that price measurement critically affects empirical inference. Using three alternative price measures---price as percentage of GNI per capita (income-relative), PPP-adjusted prices, and nominal USD prices---we show only income-relative prices yield consistent, statistically significant elasticity estimates. PPP-adjusted prices, commonly used in cross-country comparisons, produce estimates significant in just \PctPPPSigEaP\% of EaP specifications (at p$<$0.05). This methodological contribution has implications beyond broadband demand, highlighting the importance of price measurement in technology goods research where standard PPP baskets may not apply \citep{schreyer2002oecd}.

\textbf{Third}, we provide robust evidence of regional heterogeneity in price sensitivity. EaP elasticity was \EaPEURatio{} times larger than EU elasticity during 2010--2019, consistent with income effects being stronger in developing markets \citep{roller2001telecommunications}. This regional gap diminished during COVID-19, with both regions converging to near-zero elasticity---constituting empirical evidence that demand-side digital harmonisation between EaP and EU markets is not merely a policy aspiration but an observable outcome.

\textbf{Fourth}, we offer policy-relevant insights on the effectiveness of affordability interventions versus infrastructure investment. As broadband transitioned to necessity status, price-based interventions lost effectiveness while infrastructure investment and quality standards became more important. This has implications for digital inclusion policies worldwide \citep{guermazi2021digital}.

\subsection{Preview of Findings}

Our analysis yields three main findings. Pre-COVID elasticities differ sharply by region, with EaP countries \EaPEURatio{} times more price-sensitive than EU countries---a result robust across \NPerMeasure{} control specifications. Both regions then converge to near-zero elasticity during 2020--2024. Crucially, placebo tests show this decline began around 2015, pointing to secular digital integration rather than a COVID-19 shock.

The remainder of this paper proceeds as follows. Section~\ref{sec:literature} reviews relevant literature on telecommunications demand, digital transformation, and panel data methods. Section~\ref{sec:data} describes our data sources and sample construction. Section~\ref{sec:methodology} presents our empirical strategy. Section~\ref{sec:results} reports baseline and robustness results. Section~\ref{sec:discussion} discusses policy implications and limitations. Section~\ref{sec:conclusion} concludes.
